% Analytic phase plot; not deployment measurements.
% Provenance: pinned-joint theorem and b1_frontier.py worked values.
% Boundary/polygon regression: scripts/harbor-research/check_rate_regime.py.
%
% Figure brief
% Stable ID: IV/fig:rate-regime (fig:rate-regime)
% Reader question: for tolerated miss mass delta and flag budget f, is the
%   point infeasible, positively priced, or already zero-rate?
% Claim: the positive-rate wedge is bounded by f=p-delta and
%   f=1-delta/p; the zero-rate boundary is shown only where it lies inside
%   the panel, with its continuation above the panel explicitly marked.
% Relation: quantitative threshold/phase boundary.
% Required evidence: infeasible strip, positive-rate wedge, in-frame portion
%   of the zero-rate boundary, continuation cue, and the three worked values.
% Counter-reading: the top of the panel is not itself the zero-rate boundary;
%   for delta<.03 the true boundary is above f=.4.
% Visual grammar: source-owned phase plot; a single dashed continuation cue is
%   used outside the panel, never as data.
% Rejected alternative: min(1-delta/p,.4) as one plotted curve, because its
%   horizontal clipping creates a false threshold at f=.4.

\begin{figure}[H]
\centering
\pdfigurehue{pdteal}%
\begin{tikzpicture}[pd figure]
\begin{axis}[pd axis,
  width=7.8cm,height=5.6cm,
  xmin=0,xmax=0.05,ymin=0,ymax=0.4,
  xtick={0,0.01,0.02,0.03,0.04,0.05},
  ytick={0,0.05,0.10,0.20,0.30,0.40},
  xlabel={tolerated miss mass $\delta$},
  ylabel={flag budget $f$},
  axis on top,clip=false]
  % p=.05. The infeasibility floor is f=p-delta.
  % Explicit closed region polygons. Their x-domains and vertices match the
  % theorem geometry; draw=none prevents fill edges from becoming fake data.
  \path[fill=pdfocus!24,draw=none]
    (axis cs:0,0.05) -- (axis cs:0,0.4) --
    (axis cs:0.03,0.4) -- (axis cs:0.03,0.02) -- cycle;
  \path[fill=pdfocus!24,draw=none]
    (axis cs:0.03,0.02) -- (axis cs:0.03,0.4) --
    (axis cs:0.05,0) -- cycle;
  \path[fill=pdinkmuted!7,draw=none]
    (axis cs:0.03,0.4) -- (axis cs:0.05,0.4) --
    (axis cs:0.05,0) -- cycle;
  \path[fill=pdinkmuted!12,draw=none]
    (axis cs:0,0) -- (axis cs:0,0.05) --
    (axis cs:0.05,0) -- cycle;

  % Boundary strokes. The zero-rate boundary is drawn only where it is in
  % frame; its short dashed continuation follows the true slope above-left.
  \addplot[pd rule,line width=.9pt,domain=0:0.05,samples=2] {0.05-x};
  \addplot[pd focus rule,line width=1.05pt,domain=0.03:0.05,samples=60]
    {1-x/0.05};
  \draw[pd focus rule,dash pattern=on 2pt off 1.8pt,->]
    (axis cs:0.03,0.4) -- (axis cs:0.028,0.44);
  \node[pd focus label,anchor=south west]
    at (axis cs:0.005,0.405) {continues above};

  \node[pd label,anchor=west] at (axis cs:0.001,0.022) {infeasible};
  \node[pd label,anchor=west] at (axis cs:0.006,0.28) {$R(\delta,f)>0$};
  \node[pd focus label,anchor=north west] at (axis cs:0.037,0.35) {rate $=0$};

  % The three worked values, retained from the active source.
  \node[pd focus datum] at (axis cs:0,0.05) {};
  \node[pd focus label,anchor=south west] at (axis cs:0.0015,0.075) {$H(p){=}0.2864$};
  \node[pd datum] at (axis cs:0,0.10) {};
  \node[pd label,anchor=south west] at (axis cs:0.0015,0.115) {$R{=}0.1864$};
  \node[pd datum] at (axis cs:0.04,0.06) {};
  \draw[pd guide] (axis cs:0.04,0.06) -- (axis cs:0.033,0.10);
  \node[pd label,anchor=south west] at (axis cs:0.021,0.095) {$R{=}0.0087$};
\end{axis}
\end{tikzpicture}
\caption{At $p=.05$, the positive-rate wedge lies between
$f=p-\delta$ and $f=1-\delta/p$; the zero-rate boundary enters at
$(.03,.4)$ and continues above-left outside the panel.}
\label{fig:rate-regime}
\end{figure}
